\chapter{诗篇总结}
\section{诗篇的总结}
\subsection{诗篇的分类}
\begin{itemize}
\itemsep-.45em 
\item A.  赞美诗: 33, 47, 48, 68, 100, 103-105, 111, 113-118, 135, 145-150;
\item B.  感恩诗: 16, 30, 34, 40, 65-68, 75, 92, 100, 105, 107, 118, 124, 136, 146;
\item C.  君尊诗: 2, 18, 20, 21, 45, 72, 89, 101, 110, 132, 144;
\item D.  神治诗: 47, 93, 96, 97, 98 99;
\item E.  锡安诗: 48, 76, 84, 122, 134;
\item F.  祈祷诗: 3, 5, 9, 10, 12-14, 25-28, 41-44,52-57, 60, 74, 79, 80, 85, 90, 94, 140-143;
\item G.  咒诅诗: 7, 35, 58, 59, 69, 83, 109, 137, 139;
\item H.  忏悔诗: 6, 32, 38, 51, 102, 130;
\item I.  自然诗: 8, 19, 29, 65, 104, 147, 148;
\item J.  信靠诗: 4, 11, 16, 23, 27, 46, 62, 63, 91, 115, 121, 125, 131;
\item K.  智慧诗: 1, 15, 36, 37, 49, 73, 112, 127, 128, 133;
\item L.  律法诗: 19, 119;
\item M.  预言诗: 50, 78, 81, 82;
\item N.  敬拜诗: 21, 30, 33, 47, 48, 63, 65-7, 84, 95, 96, 98, 100-1, 111-3, 116, 117, 133-4, 138, 144, 148, 150;140, 143;
\item O. 安慰诗: 37, 42-3, 46-7, 91, 94, 97, 116;
\item P.  患难诗: 4, 5, 11, 28, 41, 55, 59, 64, 109, 120, 140, 143;
\item Q.  题言诗: 3, 7, 18, 30, 34, 51, 52, 54, 56, 57, 59, 60, 63;
\item R.  历史诗: 78, 81, 83, 103-6, 114, 135-7;
\item S.  训悔诗: 32, 42, 44, 45, 52-5, 74, 78, 88-9, 142;
\item T.  上行诗: 120-134;
\item U.  哀歌诗: 3-5, 7, 12-3, 26, 28, 44, 60, 74, 79, 80, 83, 85, 90, 123, 137, 142;
\item V.  乐歌诗: 4-9, 11-14;
\item W.  字母诗: 9, 10, 25, 34, 37, 111, 112, 119, 145;
\item X.  弥赛亚诗: 2, 8, 22-4, 40-1, 45, 68-9, 72, 87, 89, 102, 109, 110, 118, 132. 
\end{itemize}


\subsubsection{字母詩}
\begin{itemize}
\itemsep-.45em 
\item 字母詩定义\\
字母诗是希伯来文里面的一种体裁，它的特色是：每一节的头一个字母，是一个希伯来文的字母。每句首字的单词编排基本按照字母的顺序。希伯来文一共有22个字母，所以字母诗通常都有22节 。
\item 字母诗的目的\\
这种体裁的目的是要加深读者的印象，以致容易背诵和记忆。字母诗的难度在于要顾及字母的顺序，内容有没有连贯性，以及合不合逻辑思维等等。
\item 圣经中的字母诗
\begin{itemize}
\item 	诗篇：9.10，25，34，37，111，112，119和145篇
\item 	箴言：31：10-31
\item 	耶利米哀歌：1，2，3，和4
\item 	那鸿书：1：2-10
\end{itemize}

\item 注明
\begin{itemize}
\item 	诗篇》第9篇和第10篇在七十士译本内，原来是合在一起的，是同一首完整的字母诗 。第9篇用前11个字母，而第10篇用后11个字母 , 。
\item 	《诗篇》第111篇和112 篇，在中文翻译的圣经和英文翻译的圣经都只有10节经文；但是在希伯来原文的旧约圣经里面，这两篇诗都有22句诗句，正好就用到了希伯来文的22个字母，可惜翻译成别国文字的时候，失去了字母诗的特色和韵味。
\item 	《诗篇》第119篇也是字母诗，它有176节，它的节数是22个字母的倍数，等于每8节用同一个字母，共有22段 。
\end{itemize}


\end{itemize}

\subsubsection{詩篇中的诗章Psalms（赞美天父的创造类）}
靈交詩-讚美類（8，20，33，34，40，47，71，92，95）
\begin{itemize}
\item I.	時刻都當讚美神(個人獨處，集體中讚美) 
\begin{itemize}
\item A.	因神是至高者（诗47）因为耶和华是至高者，治理全地。他治理万国，坐在他的圣宝座上。並且神超乎是大王，超乎萬神之上（詩95:3）。
\item B.	神偉大的創造（詩8:3），最高的山峰,最低的地，海洋,旱地都屬於他（詩95:4-5）
\item C.	安息日的時候歌頌耶和華至高的名（詩92小標題）
\item D.	遭難的日子
\begin{itemize}
\item 1.	大衛躲避掃羅，逃到迦特王那裡。因怕迦特王嫉妒自己的能力，裝瘋賣傻，從而糟迦特王厭棄，而逃脫（詩34小標題,撒上21:10-15）。
\item 2.	從禍患和淤泥裡，把大衛拉上來，並且使大衛腳步穩當（詩40:2）。
\item 3.	自己的罪孽太多，使自己心寒膽戰的時候（诗40:12）。
\item 4.	仇敵追趕，並認為神不再搭救的時候，作者卻要越發讚美（詩71:10，14）。
\item 5.	神是我们的坚固避难所（诗71:8）
\end{itemize}
\item E.	重大事情發生前，獻祭倚靠神 （詩20）。大衛在打仗前，獻祭，祈求耶和華的拯救和同在。
\item F.	記念耶和華奇妙的作為，回顾创世纪和出埃及记的时候，帶領選民到立約之地，（詩105）。
\item G.	個人和集體讚美。 
\begin{itemize}
\item 個人讚美：詩篇8，40，71，92
\item 集體讚美：詩篇20，33，34，47，95
\end{itemize}
\end{itemize}
\item II.	發出讚美提升我們的靈性
\begin{itemize}
\item A.	感恩神的眷顧（詩8:4）。創造天地偉大的神，竟然眷顧這微小的我們讓我們來管理神創造的萬物。
\item B.	我們的心會歡喜，因為我們依靠耶和華（詩33:21）。也因為依靠耶和華是有福的（詩40:4）。歌頌耶和華可以使我們高興歡呼，歡喜（詩92:4,詩105:3-4）。
\item C.	知道神會救護我們（詩20:6）。在遭難的日子，祈求和讚美，會驅除內心的恐懼，反而領悟到神會救護我們，心裡就會有盼望。雖然在困苦窮乏中，稱頌耶和華的名為大，但相信神依然顧念我（詩40:17）。
\item D.	仰望耶和華必不蒙羞（詩34:5）。一個曾經風光一時的大衛王，此時要裝瘋賣傻，但是因為讚美就對神有了仰望，內心就不會蒙羞。
\item E.	感恩珍惜自己上帝儿女的身份
\begin{itemize}
\item 1.	我们是耶和华草场的羊，是耶和华手下的民（詩95:7）。
\item 2.	耶和华高举了我们的角，我们是被新油膏的（诗92:10）。
\end{itemize}
\end{itemize}
\item III.	神配得我们的赞美
\begin{itemize}
\item A.	耶和華的創造，偉大，和他的本性
\begin{itemize}
\item 1.	耶和華言語正直，並充滿了慈愛（詩33:4-6）
\item 2.	神偉大的創造。整個世界，包括巨大的海洋都是藉著神而創造的（詩33:6-11）。地的深處還是最高的高峰，無論是海洋還是旱地都是他手所造的（詩95:4-5）。
\item 3.	耶和華是正直的（詩92:15）。
\item 4.	耶和华是至高的神，直到永远（诗92:8）。
\item 5.	耶和華為大王，大神，超乎萬神之上（詩95:3）。
\end{itemize}
\item B.	耶和華的籌算堅立和心思及深
\begin{itemize}
\item 1.	耶和華的籌算永遠立定（詩33:10-11）。儘管各國，各民都有自己的打算，但若不是耶和華的籌算，都會變為烏有。
\item 2.	耶和華的心思及其深，愚妄人是不能明白的（詩92:5-6）。
\end{itemize}
\item C.	耶和華的美善和賜福
\begin{itemize}
\item 1.	投靠耶和華的人是有福的，敬畏和尋求耶和華的一無所缺（詩34:8-9）。
\item 2.	我們可以按照神的命令，舌頭不出惡言，行善，尋求和睦，就能得享美福（詩34:12-14）。
\item 3.	義人發旺，栽于耶和華的殿中，到年老的時候仍然滿了汁漿而常发青（詩92:13-14）
\end{itemize}
\item D.	困苦患難呼求的時候，應允，看顧我們的神拯救我們的神
\begin{itemize}
\item 1.	願神应允，高舉，堅固，成就我們（詩20：1-5）
\item 2.	耶和華是我們的依靠，打仗有人靠車，有人靠馬，我們依靠耶和華神的名（詩20:7）
\item 3.	耶和華在敬畏他的人四圍安營扎寨搭救他們，從而保守大衛脫離了患難（詩34:6-7）
\item 4.	耶和華話靠近傷心的人，拯救靈性痛悔的人（詩34:18）。
\item 5.	雖然經歷苦難，但是神保守我們，竟然一根骨頭都不能折斷（詩34:19-20）。
\item 6.	耶和華的的看顧和救護（詩33:12-19）。我們的神雖然高高在天，但是他密切關注我們，我們在他面前就是透明人一樣。同時，耶和華又會看顧敬畏他的人，保守敬畏他的人脫離死亡，在饑荒中也能存活。
\end{itemize}
\item E.	耶和華的保護和救恩
1.	耶和華是我們的磐石（詩40:2，71:3，92:15，95:1 ）
2.	耶和華救贖我們的靈魂，凡投靠他的必不至定罪（詩34:22）。
3.	神給了我們救恩（詩20:5，40:10，和71:15）。
\item F.	公義的神。先祖們在曠野試探神40年，最終神厭煩了那個世代，並且在怒中起誓不可進入神的安息（詩95:8-11）.
\item G.	守約的神（詩105:8-11,42-43）。神與亞伯拉罕立約，必將迦南地賜給他們，以色列人在埃及為奴400年之後，神最終帶領選民進入了所賜之地。
\item H.	行奇事的神（诗105）.
\begin{itemize}
\item 1.	神保护，不容人欺负他的子民（诗105:12-15）。
亚伯拉罕在基拉耳寄居的时候，认为这地方的人不惧怕神，会因为他妻子的原因杀他，就让撒拉称他为哥哥。基拉耳王亚比米勒将撒拉取来，夜间耶和华对亚比米勒说话。不仅归还了亚比米勒，还称亚伯拉罕为先知，让亚伯拉罕为亚比米勒祷告，并且借着祷告，医治了他妻子和众女仆，使他们怀孕（创世纪20章）。
\item 2.	神預備約瑟去埃及做準備，並成為王家之主（105:17-23）
创世纪37章，约瑟被哥哥卖到埃及，在39章，因为波提乏的妻子陷害他，坐监。虽然在狱中帮助酒政，酒政出去后忘记了他 （40章），一直又过了2年，酒政才想起来约瑟可以为法老解梦，最终法拉派他治理埃及全地（41章）。
\item 3.	神使百姓在埃及生養眾多（105:24）。
以色列的众子到埃及的时候共有70人（出1：5）。之后以色列人生养众多，出埃及的时候，除了妇人和孩子，步行的男人约有60万（出13:37）。
\item 4.	神打发他的僕人摩西和亞倫在敵人中行了很多神蹟（105:26-38）。
当摩西和亚伦带领百姓出埃及的时候，行了亚伦的杖变蛇，水变血之灾，蛙灾，虱灾，蝇灾，畜疫灾，疮灾，雹灾，蝗灾，黑暗之灾，杀长子之灾，最终带领百姓出了埃及（创7-11章）。
\item 5.	帶領百姓而出，雲彩遮蓋，夜間有火光照，時刻同在，並且供應生活所需的食物和水源（105:39-41）。
\end{itemize}
\item I.	看顾我们一生的神。
\begin{itemize}
\item 1.	从出母胎的时候被神扶持（诗71:6）。
\item 2.	年幼的时候，神是我们的依靠（诗71:5）。
\item 3.	年老的时候，力气衰弱的时候，不要离弃我（诗71：9）。
\item 4.	年老的时候依然栽于耶和华的殿中，发旺在神的院里（诗92:13-14）。
\end{itemize}
\item J.	赐给我们荣耀的神
\begin{itemize}
\item 1.	神使万民服在我们以下，叫列邦服在我们脚下（诗71:2-3）。
\item 2.	神使我们管理他手所造的牛羊，兽，鸟和鱼（诗8:6-8）。
\end{itemize}
\end{itemize}
\item IV.	詩人怎麼讚美神？
\begin{itemize}
\item A.	呼求神（詩8），求告神的名（诗105:1）。
\item B.	诗歌弹琴
\begin{itemize}
\item 1.	彈琴唱歌讚美神（詩33）。
\item 2.	用十弦的乐器和瑟，用琴弹优雅的声音（詩92:2-3）。
\item 3.	拍掌欢呼（诗47:1）。我们现在敬拜的时候有时也是边唱边拍掌。
\end{itemize}
\item C.	感恩称谢
\begin{itemize}
\item 1. 称谢耶和华，歌颂神至高的名（诗92:1）。
\item 2. 感謝神，用詩歌向他歡呼（诗95:2）
\end{itemize}
\item D.	時時稱頌耶和華（詩34:1），早晨傳揚神的慈愛，每夜傳揚神的信實（詩92:2-3），在万民中传扬神奇妙的作为（诗105:1-2）。
\item E.	屈身敬拜，在我們耶和華面前跪下（詩95:6）。
\item F.	用悟性歌颂（诗47:7）。我们歌颂的时候不仅仅是表面的言语，还要搭配内化，领悟，祷告式的赞美和歌颂。
\end{itemize}
\end{itemize}


\subsubsection{詩篇中的颂词Hymns（个人生命中经历子的救赎）}

\subsubsection{詩篇中的灵歌Spiritual Songs（对圣灵的渴慕类）}

\subsection{诗篇中的神学}
\begin{itemize}
\itemsep-.45em  
\item 神本体论
\item 基督论
\item 圣灵论
\item 天使论
\item 人论
\item 罪论
\item 救赎论
\item 末世论
\end{itemize}

\subsubsection{诗篇中的基督}
\begin{itemize}
\item 基督
\begin{itemize}
\itemsep-.45em  
\item A.  基督的国度: 诗篇2
\item B.  基督的人性: 诗篇8
\item C.  基督的复活: 诗篇16
\item D.  基督的受难: 诗篇22
\item E.  基督是牧者: 诗篇23
\item F.  基督大君王: 诗篇24
\item G.  基督的顺服: 诗篇40
\item H.  基督被出卖: 诗篇41
\item I.  基督的婚娶: 诗篇45 (诗篇中的雅歌)
\item J.  基督的国度: 诗篇68，72,89（大卫之约）
\item K.  基督的受苦: 诗篇69 (客西马尼诗)
\item L.  基督的被弃: 诗篇102; 109
\item M.  基督为祭司: 诗篇110 (麦基洗德诗)
\item N.  基督的被弃与高升: 诗篇118
\item O.  基督的宝座: 诗篇132 (宝座诗) 
\end{itemize}
\item 基督的生平
\begin{itemize}
\itemsep-.45em  
\item A.  神性: 诗45:6, 11; 来1:8
\item B.  永存: 诗45:17; 61:6-7; 72:17; 102:25-27; 来1:10
\item C.  权能: 诗72:8; 103:19; 启19:16
\item D.  神子: 诗2:7; 110:1; 来1:5; 太22:40-45;
\item E.  降生: 诗40:6-8; 来10;5-7
\item F.  圣洁: 诗45:7; 89:18-19; 来1:9
\item G.  热心: 诗69:9; 约2:17
\item H.  先知: 诗22:22; 40:9-10
\item I.  祭司: 诗110; 来5:6;
\item J.  受苦: 诗69:9; 罗15:3
\item K.  被卖: 诗41:9; 路22:48
\item L.  受死: 诗22:1-21; 31:5;
\item M.  被弃: 诗118:22-23; 太21:42
\item N.  复活: 诗2:7; 16:10; 徒2:27; 13:35
\item O.  升天: 诗68:18; 弗4:8
\item P.  再来: 诗96-98; 帖后1:7-9
\item Q. 为王: 诗2:6; 89:18-19; 徒5:31; 启19:16
\item R.  得胜: 诗110:5-6; 启6:17
\end{itemize}
\end{itemize}

\subsubsection{诗篇中的祷告}

\subsubsection{诗篇中的预言}


